\iftestbuild % TEST-ONLY-BEGIN
\setcounter{section}{3}\setcounter{theorem}{43}
\fi % TEST-ONLY-END
% Source: book pp. 258--269 (PDF 272--283); no solutions.
% \axeqno: see 7C.tex
\DeclareRobustCommand\axeqno{\refstepcounter{theorem}\tag*{\thetheorem}}
\section{Isometries, Unitary Operators, and Matrix Factorization}

\subsection{Isometries}

Linear maps that preserve norms are sufficiently important to deserve a name.

\begin{definition}[Isometry]\label{ax:7.44}
A linear map $S\in\Lin(V,W)$ is called an \emph{isometry} if
\[ \norm{Sv}=\norm{v} \]
for every $v\in V$. In other words, a linear map is an isometry if it preserves
norms.
\end{definition}

\marginnote{The Greek word \textbf{isos} means equal; the Greek word
\textbf{metron} means measure. Thus \textbf{isometry} literally means equal
measure.}%
If $S\in\Lin(V,W)$ is an isometry and $v\in V$ is such that $Sv=0$, then
\[ \norm{v}=\norm{Sv}=\norm{0}=0, \]
which implies that $v=0$. Thus every isometry is injective.

\begin{example}[Orthonormal basis maps to orthonormal list $\implies$ isometry]%
\label{ax:7.45}
Suppose $e_1,\dots,e_n$ is an orthonormal basis of $V$ and $g_1,\dots,g_n$ is an
orthonormal list in $W$. Let $S\in\Lin(V,W)$ be the linear map such that
$Se_k=g_k$ for each $k=1,\dots,n$. To show that $S$ is an isometry, suppose
$v\in V$. Then
\[ v=\ip{v}{e_1}e_1+\cdots+\ip{v}{e_n}e_n \axeqno\label{ax:7.46} \]
and
\[ \norm{v}^2=\abs{\ip{v}{e_1}}^2+\cdots+\abs{\ip{v}{e_n}}^2,
   \axeqno\label{ax:7.47} \]
where we have used \ref{ax:6.30}(b). Applying $S$ to both sides of
\ref{ax:7.46} gives
\[ Sv=\ip{v}{e_1}Se_1+\cdots+\ip{v}{e_n}Se_n
     =\ip{v}{e_1}g_1+\cdots+\ip{v}{e_n}g_n. \]
Thus
\[ \norm{Sv}^2=\abs{\ip{v}{e_1}}^2+\cdots+\abs{\ip{v}{e_n}}^2.
   \axeqno\label{ax:7.48} \]
Comparing \ref{ax:7.47} and \ref{ax:7.48} shows that $\norm{v}=\norm{Sv}$. Thus
$S$ is an isometry.
\end{example}

The next result gives conditions equivalent to being an isometry. The
equivalence of (a) and (c) shows that a linear map is an isometry if and only if
it preserves inner products. The equivalence of (a) and (d) shows that a linear
map is an isometry if and only if it maps some orthonormal basis to an
orthonormal list. Thus the isometries given by Example~\ref{ax:7.45} include all
isometries. Furthermore, a linear map is an isometry if and only if it maps
every orthonormal basis to an orthonormal list [because whether or not (a) holds
does not depend on the basis $e_1,\dots,e_n$].

The equivalence of (a) and (e) in the next result shows that a linear map is an
isometry if and only if the columns of its matrix (with respect to any
orthonormal bases) form an orthonormal list. Here we are identifying the columns
of an $m$-by-$n$ matrix with elements of $\F^m$ and then using the Euclidean
inner product on $\F^m$.

\begin{result}[Characterizations of isometries]\label{ax:7.49}
Suppose $S\in\Lin(V,W)$. Suppose $e_1,\dots,e_n$ is an orthonormal basis of $V$
and $f_1,\dots,f_m$ is an orthonormal basis of $W$. Then the following are
equivalent.
\begin{parts}
\item $S$ is an isometry.
\item $S^*S=I$.
\item $\ip{Su}{Sv}=\ip{u}{v}$ for all $u,v\in V$.
\item $Se_1,\dots,Se_n$ is an orthonormal list in $W$.
\item The columns of $\Mat\bigl(S,(e_1,\dots,e_n),(f_1,\dots,f_m)\bigr)$ form an
  orthonormal list in $\F^m$ with respect to the Euclidean inner product.
\end{parts}
\end{result}

\begin{proof}
First suppose (a) holds, so $S$ is an isometry. If $v\in V$ then
\[ \bigl\langle(I-S^*S)v,v\bigr\rangle=\ip{v}{v}-\ip{S^*Sv}{v}
   =\norm{v}^2-\ip{Sv}{Sv}=\norm{v}^2-\norm{Sv}^2=0. \]
Hence the self-adjoint operator $I-S^*S$ equals $0$ (by \ref{ax:7.16}). Thus
$S^*S=I$, proving that (a) implies (b).

Now suppose (b) holds, so $S^*S=I$. If $u,v\in V$ then
\[ \ip{Su}{Sv}=\ip{S^*Su}{v}=\ip{Iu}{v}=\ip{u}{v}, \]
proving that (b) implies (c).

Now suppose that (c) holds, so $\ip{Su}{Sv}=\ip{u}{v}$ for all $u,v\in V$. Thus
if $j,k\in\{1,\dots,n\}$, then
\[ \ip{Se_j}{Se_k}=\ip{e_j}{e_k}. \]
Hence $Se_1,\dots,Se_n$ is an orthonormal list in $W$, proving that (c) implies
(d).

Now suppose that (d) holds, so $Se_1,\dots,Se_n$ is an orthonormal list in $W$.
Let $A=\Mat\bigl(S,(e_1,\dots,e_n),(f_1,\dots,f_m)\bigr)$. If
$k,r\in\{1,\dots,n\}$, then
\[ \sum_{j=1}^mA_{j,k}\overline{A_{j,r}}
   =\Bigl\langle\sum_{j=1}^mA_{j,k}f_j,\sum_{j=1}^mA_{j,r}f_j\Bigr\rangle
   =\ip{Se_k}{Se_r}
   =\begin{cases}1&\text{if }k=r,\\0&\text{if }k\ne r.\end{cases}
   \axeqno\label{ax:7.50} \]
The left side of \ref{ax:7.50} is the inner product in $\F^m$ of columns $k$ and
$r$ of $A$. Thus the columns of $A$ form an orthonormal list in $\F^m$, proving
that (d) implies (e).

Now suppose (e) holds, so the columns of the matrix $A$ defined in the paragraph
above form an orthonormal list in $\F^m$. Then \ref{ax:7.50} shows that
$Se_1,\dots,Se_n$ is an orthonormal list in $W$. Thus Example~\ref{ax:7.45},
with $Se_1,\dots,Se_n$ playing the role of $g_1,\dots,g_n$, shows that $S$ is an
isometry, proving that (e) implies (a).
\end{proof}

\marginnote[-2em]{Although the words ``unitary'' and ``isometry'' mean the same thing
for operators on finite-dimensional inner product spaces, remember that a
unitary operator maps a vector space to itself, while an isometry maps a vector
space to another (possibly different) vector space.}%
As previously noted, every isometry is injective. Every injective operator on a
finite-dimensional vector space is invertible (see \ref{ax:3.65}). A standing
assumption for this chapter is that $V$ is a finite-dimensional inner product
space. Thus we could delete the word ``invertible'' from the definition above
without changing the meaning. The unnecessary word ``invertible'' has been
retained in the definition above for consistency with the definition readers
may encounter when learning about inner product spaces that are not necessarily
finite-dimensional.


See Exercises 1 and 11 for additional conditions that are equivalent to being
an isometry.

\subsection{Unitary Operators}

In this subsection, we confine our attention to linear maps from a vector space
to itself. In other words, we will be working with operators.

\begin{definition}[Unitary operator]\label{ax:7.51}
An operator $S\in\Lin(V)$ is called \emph{unitary} if $S$ is an invertible
isometry.
\end{definition}


\begin{example}[Rotation of $\R^2$]\label{ax:7.52}
Suppose $\theta\in\R$ and $S$ is the operator on $\F^2$ whose matrix with respect
to the standard basis of $\F^2$ is
\[ \begin{pmatrix}\cos\theta&-\sin\theta\\\sin\theta&\cos\theta\end{pmatrix}. \]
The two columns of this matrix form an orthonormal list in $\F^2$; hence $S$ is
an isometry [by the equivalence of (a) and (e) in \ref{ax:7.49}]. Thus $S$ is a
unitary operator.

If $\F=\R$, then $S$ is the operator of counterclockwise rotation by $\theta$
radians around the origin of $\R^2$. This observation gives us another way to
think about why $S$ is an isometry, because each rotation around the origin of
$\R^2$ preserves norms.
\end{example}

The next result (\ref{ax:7.53}) lists several conditions that are equivalent to
being a unitary operator. All the conditions equivalent to being an isometry in
\ref{ax:7.49} should be added to this list. The extra conditions in
\ref{ax:7.53} arise because of limiting the context to linear maps from a vector
space to itself. For example, \ref{ax:7.49} shows that a linear map
$S\in\Lin(V,W)$ is an isometry if and only if $S^*S=I$, while \ref{ax:7.53}
shows that an operator $S\in\Lin(V)$ is a unitary operator if and only if
$S^*S=SS^*=I$.

Another difference is that \ref{ax:7.49}(d) mentions an orthonormal list, while
\ref{ax:7.53}(d) mentions an orthonormal basis. Also, \ref{ax:7.49}(e) mentions
the columns of $\Mat(T)$, while \ref{ax:7.53}(e) mentions the rows of $\Mat(T)$.
Furthermore, $\Mat(T)$ in \ref{ax:7.49}(e) is with respect to an orthonormal
basis of $V$ and an orthonormal basis of $W$, while $\Mat(T)$ in
\ref{ax:7.53}(e) is with respect to a single basis of $V$ doing double duty.

\begin{result}[Characterizations of unitary operators]\label{ax:7.53}
Suppose $S\in\Lin(V)$. Suppose $e_1,\dots,e_n$ is an orthonormal basis of $V$.
Then the following are equivalent.
\begin{parts}
\item $S$ is a unitary operator.
\item $S^*S=SS^*=I$.
\item $S$ is invertible and $S^{-1}=S^*$.
\item $Se_1,\dots,Se_n$ is an orthonormal basis of $V$.
\item The rows of $\Mat\bigl(S,(e_1,\dots,e_n)\bigr)$ form an orthonormal basis
  of $\F^n$ with respect to the Euclidean inner product.
\item $S^*$ is a unitary operator.
\end{parts}
\end{result}

\begin{proof}
First suppose (a) holds, so $S$ is a unitary operator. Hence
\[ S^*S=I \]
by the equivalence of (a) and (b) in \ref{ax:7.49}. Multiply both sides of this
equation by $S^{-1}$ on the right, getting $S^*=S^{-1}$. Thus
$SS^*=SS^{-1}=I$, as desired, proving that (a) implies (b).

The definitions of invertible and inverse show that (b) implies (c).

Now suppose (c) holds, so $S$ is invertible and $S^{-1}=S^*$. Thus $S^*S=I$.
Hence $Se_1,\dots,Se_n$ is an orthonormal list in $V$, by the equivalence of (b)
and (d) in \ref{ax:7.49}. The length of this list equals $\dim V$. Thus
$Se_1,\dots,Se_n$ is an orthonormal basis of $V$, proving that (c) implies (d).

Now suppose (d) holds, so $Se_1,\dots,Se_n$ is an orthonormal basis of $V$. The
equivalence of (a) and (d) in \ref{ax:7.49} shows that $S$ is a unitary
operator. Thus
\[ (S^*)^*S^*=SS^*=I, \]
where the last equation holds because we already showed that (a) implies (b) in
this result. The equation above and the equivalence of (a) and (b) in
\ref{ax:7.49} show that $S^*$ is an isometry. Thus the columns of
$\Mat\bigl(S^*,(e_1,\dots,e_n)\bigr)$ form an orthonormal basis of $\F^n$ [by the
equivalence of (a) and (e) of \ref{ax:7.49}]. The rows of
$\Mat\bigl(S,(e_1,\dots,e_n)\bigr)$ are the complex conjugates of the columns of
$\Mat\bigl(S^*,(e_1,\dots,e_n)\bigr)$. Thus the rows of
$\Mat\bigl(S,(e_1,\dots,e_n)\bigr)$ form an orthonormal basis of $\F^n$, proving
that (d) implies (e).

Now suppose (e) holds. Thus the columns of $\Mat\bigl(S^*,(e_1,\dots,e_n)\bigr)$
form an orthonormal basis of $\F^n$. The equivalence of (a) and (e) in
\ref{ax:7.49} shows that $S^*$ is an isometry, proving that (e) implies (f).

Now suppose (f) holds, so $S^*$ is a unitary operator. The chain of implications
we have already proved in this result shows that (a) implies (f). Applying this
result to $S^*$ shows that $(S^*)^*$ is a unitary operator, proving that (f)
implies (a).

We have shown that (a) $\Rightarrow$ (b) $\Rightarrow$ (c) $\Rightarrow$ (d)
$\Rightarrow$ (e) $\Rightarrow$ (f) $\Rightarrow$ (a), completing the proof.
\end{proof}

Recall our analogy between $\C$ and $\Lin(V)$. Under this analogy, a complex
number $z$ corresponds to an operator $S\in\Lin(V)$, and $\overline{z}$
corresponds to $S^*$. The real numbers $(z=\overline{z})$ correspond to the
self-adjoint operators $(S=S^*)$, and the nonnegative numbers correspond to the
(badly named) positive operators.

Another distinguished subset of $\C$ is the unit circle, which consists of the
complex numbers $z$ such that $\abs{z}=1$. The condition $\abs{z}=1$ is
equivalent to the condition $\overline{z}z=1$. Under our analogy, this
corresponds to the condition $S^*S=I$, which is equivalent to $S$ being a
unitary operator. Hence the analogy shows that the unit circle in $\C$
corresponds to the set of unitary operators. In the next two results, this
analogy appears in the eigenvalues of unitary operators. Also see Exercise~15
for another example of this analogy.

\begin{result}[Eigenvalues of unitary operators have absolute value $1$]%
\label{ax:7.54}
Suppose $\lambda$ is an eigenvalue of a unitary operator. Then $\abs{\lambda}=1$.
\end{result}

\begin{proof}
Suppose $S\in\Lin(V)$ is a unitary operator and $\lambda$ is an eigenvalue of
$S$. Let $v\in V$ be such that $v\ne0$ and $Sv=\lambda v$. Then
\[ \abs{\lambda}\,\norm{v}=\norm{\lambda v}=\norm{Sv}=\norm{v}. \]
Thus $\abs{\lambda}=1$, as desired.
\end{proof}

The next result characterizes unitary operators on finite-dimensional complex
inner product spaces, using the complex spectral theorem as the main tool.

\begin{result}[Description of unitary operators on complex inner product spaces]%
\label{ax:7.55}
Suppose $\F=\C$ and $S\in\Lin(V)$. Then the following are equivalent.
\begin{parts}
\item $S$ is a unitary operator.
\item There is an orthonormal basis of $V$ consisting of eigenvectors of $S$
  whose corresponding eigenvalues all have absolute value $1$.
\end{parts}
\end{result}

\begin{proof}
Suppose (a) holds, so $S$ is a unitary operator. The equivalence of (a) and (b)
in \ref{ax:7.53} shows that $S$ is normal. Thus the complex spectral theorem
(\ref{ax:7.31}) shows that there is an orthonormal basis $e_1,\dots,e_n$ of $V$
consisting of eigenvectors of $S$. Every eigenvalue of $S$ has absolute value
$1$ (by \ref{ax:7.54}), completing the proof that (a) implies (b).

Now suppose (b) holds. Let $e_1,\dots,e_n$ be an orthonormal basis of $V$
consisting of eigenvectors of $S$ whose corresponding eigenvalues
$\lambda_1,\dots,\lambda_n$ all have absolute value $1$. Then $Se_1,\dots,Se_n$
is also an orthonormal basis of $V$ because
\[ \ip{Se_j}{Se_k}=\ip{\lambda_je_j}{\lambda_ke_k}
   =\lambda_j\overline{\lambda_k}\ip{e_j}{e_k}
   =\begin{cases}0&\text{if }j\ne k,\\1&\text{if }j=k\end{cases} \]
for all $j,k=1,\dots,n$. Thus the equivalence of (a) and (d) in \ref{ax:7.53}
shows that $S$ is unitary, proving that (b) implies (a).
\end{proof}

\subsection{QR Factorization}

In this subsection, we shift our attention from operators to matrices. This
switch should give you good practice in identifying an operator with a square
matrix (after picking a basis of the vector space on which the operator is
defined). You should also become more comfortable with translating concepts and
results back and forth between the context of operators and the context of
square matrices.

When starting with $n$-by-$n$ matrices instead of operators, unless otherwise
specified assume that the associated operators live on $\F^n$ (with the
Euclidean inner product) and that their matrices are computed with respect to
the standard basis of $\F^n$.

We begin by making the following definition, transferring the notion of a
unitary operator to a unitary matrix.

\begin{definition}[Unitary matrix]\label{ax:7.56}
An $n$-by-$n$ matrix is called \emph{unitary} if its columns form an orthonormal
list in $\F^n$.
\end{definition}

In the definition above, we could have replaced ``orthonormal list in $\F^n$''
with ``orthonormal basis of $\F^n$'' because every orthonormal list of length
$n$ in an $n$-dimensional inner product space is an orthonormal basis. If
$S\in\Lin(V)$ and $e_1,\dots,e_n$ and $f_1,\dots,f_n$ are orthonormal bases of
$V$, then $S$ is a unitary operator if and only if
$\Mat\bigl(S,(e_1,\dots,e_n),(f_1,\dots,f_n)\bigr)$ is a unitary matrix, as
shown by the equivalence of (a) and (e) in \ref{ax:7.49}. Also note that we
could also have replaced ``columns'' in the definition above with ``rows'' by
using the equivalence between conditions (a) and (e) in \ref{ax:7.53}.

The next result, whose proof will be left as an exercise for the reader, gives
some equivalent conditions for a square matrix to be unitary. In (c), $Qv$
denotes the matrix product of $Q$ and $v$, identifying elements of $\F^n$ with
$n$-by-$1$ matrices (sometimes called column vectors). The norm in (c) below is
the usual Euclidean norm on $\F^n$ that comes from the Euclidean inner product.
In (d), $Q^*$ denotes the conjugate transpose of the matrix $Q$, which
corresponds to the adjoint of the associated operator.

\begin{result}[Characterizations of unitary matrices]\label{ax:7.57}
Suppose $Q$ is an $n$-by-$n$ matrix. Then the following are equivalent.
\begin{parts}
\item $Q$ is a unitary matrix.
\item The rows of $Q$ form an orthonormal list in $\F^n$.
\item $\norm{Qv}=\norm{v}$ for every $v\in\F^n$.
\item $Q^*Q=QQ^*=I$, the $n$-by-$n$ matrix with $1$'s on the diagonal and $0$'s
  elsewhere.
\end{parts}
\end{result}

The QR factorization stated and proved below is the main tool in the widely
used QR algorithm (not discussed here) for finding good approximations to
eigenvalues and eigenvectors of square matrices. In the result below, if the
matrix $A$ is in $\F^{n,n}$, then the matrices $Q$ and $R$ are also in
$\F^{n,n}$.

\begin{result}[QR factorization]\label{ax:7.58}
Suppose $A$ is a square matrix with linearly independent columns. Then there
exist unique matrices $Q$ and $R$ such that $Q$ is unitary, $R$ is upper
triangular with only positive numbers on its diagonal, and
\[ A=QR. \]
\end{result}

\begin{proof}
Let $v_1,\dots,v_n$ denote the columns of $A$, thought of as elements of
$\F^n$. Apply the Gram--Schmidt procedure (\ref{ax:6.32}) to the list
$v_1,\dots,v_n$, getting an orthonormal basis $e_1,\dots,e_n$ of $\F^n$ such that
\[ \spn(v_1,\dots,v_k)=\spn(e_1,\dots,e_k) \axeqno\label{ax:7.59} \]
for each $k=1,\dots,n$. Let $R$ be the $n$-by-$n$ matrix defined by
\[ R_{j,k}=\ip{v_k}{e_j}, \]
where $R_{j,k}$ denotes the entry in row $j$, column $k$ of $R$. If $j>k$, then
$e_j$ is orthogonal to $\spn(e_1,\dots,e_k)$ and hence $e_j$ is orthogonal to
$v_k$ (by \ref{ax:7.59}). In other words, if $j>k$ then $\ip{v_k}{e_j}=0$. Thus
$R$ is an upper-triangular matrix.

Let $Q$ be the unitary matrix whose columns are $e_1,\dots,e_n$. If
$k\in\{1,\dots,n\}$, then the $k^{\text{th}}$ column of $QR$ equals a linear
combination of the columns of $Q$, with the coefficients for the linear
combination coming from the $k^{\text{th}}$ column of $R$---see
\ref{ax:3.51}(a). Hence the $k^{\text{th}}$ column of $QR$ equals
\[ \ip{v_k}{e_1}e_1+\cdots+\ip{v_k}{e_k}e_k, \]
which equals $v_k$ [by \ref{ax:6.30}(a)], the $k^{\text{th}}$ column of $A$.
Thus $A=QR$, as desired.

The equations defining the Gram--Schmidt procedure (see \ref{ax:6.32}) show
that each $v_k$ equals a positive multiple of $e_k$ plus a linear combination of
$e_1,\dots,e_{k-1}$. Thus each $\ip{v_k}{e_k}$ is a positive number. Hence all
entries on the diagonal of $R$ are positive numbers, as desired.

Finally, to show that $Q$ and $R$ are unique, suppose we also have
$A=\widehat{Q}\widehat{R}$, where $\widehat{Q}$ is unitary and $\widehat{R}$ is
upper triangular with only positive numbers on its diagonal. Let
$q_1,\dots,q_n$ denote the columns of $\widehat{Q}$. Thinking of matrix
multiplication as above, we see that each $v_k$ is a linear combination of
$q_1,\dots,q_k$, with the coefficients coming from the $k^{\text{th}}$ column of
$\widehat{R}$. This implies that $\spn(v_1,\dots,v_k)=\spn(q_1,\dots,q_k)$ and
$\ip{v_k}{q_k}>0$. The uniqueness of the orthonormal lists satisfying these
conditions (see Exercise~10 in Section~6B) now shows that $q_k=e_k$ for each
$k=1,\dots,n$. Hence $\widehat{Q}=Q$, which then implies that $\widehat{R}=R$,
completing the proof of uniqueness.
\end{proof}

The proof of the QR factorization shows that the columns of the unitary matrix
can be computed by applying the Gram--Schmidt procedure to the columns of the
matrix to be factored. The next example illustrates the computation of the QR
factorization based on the proof that we just completed.

\begin{example}[QR factorization of a $3$-by-$3$ matrix]\label{ax:7.60}
To find the QR factorization of the matrix
\[ A=\begin{pmatrix}1&2&1\\0&1&-4\\0&3&2\end{pmatrix}, \]
follow the proof of \ref{ax:7.58}. Thus set $v_1,v_2,v_3$ equal to the columns
of $A$:
\[ v_1=(1,0,0),\quad v_2=(2,1,3),\quad v_3=(1,-4,2). \]
Apply the Gram--Schmidt procedure to $v_1,v_2,v_3$, producing the orthonormal
list
\[ e_1=(1,0,0),\quad
   e_2=\Bigl(0,1/\sqrt{10},3/\sqrt{10}\Bigr),\quad
   e_3=\Bigl(0,-3/\sqrt{10},1/\sqrt{10}\Bigr). \]
Still following the proof of \ref{ax:7.58}, let $Q$ be the unitary matrix whose
columns are $e_1,e_2,e_3$:
\[ Q=\begin{pmatrix}1&0&0\\[2pt]
     0&1/\sqrt{10}&-3/\sqrt{10}\\[6pt]
     0&3/\sqrt{10}&1/\sqrt{10}\end{pmatrix}. \]
As in the proof of \ref{ax:7.58}, let $R$ be the $3$-by-$3$ matrix whose entry
in row $j$, column $k$ is $\ip{v_k}{e_j}$, which gives
\[ R=\begin{pmatrix}1&2&1\\[2pt]
     0&\sqrt{10}&\sqrt{10}/5\\[6pt]
     0&0&7\sqrt{10}/5\end{pmatrix}. \]
Note that $R$ is indeed an upper-triangular matrix with only positive numbers on
the diagonal, as required by the QR factorization.

Now matrix multiplication can verify that $A=QR$ is the desired factorization
of $A$:
\[ QR=\begin{pmatrix}1&0&0\\[2pt]
     0&1/\sqrt{10}&-3/\sqrt{10}\\[6pt]
     0&3/\sqrt{10}&1/\sqrt{10}\end{pmatrix}
     \begin{pmatrix}1&2&1\\[2pt]
     0&\sqrt{10}&\sqrt{10}/5\\[6pt]
     0&0&7\sqrt{10}/5\end{pmatrix}
    =\begin{pmatrix}1&2&1\\0&1&-4\\0&3&2\end{pmatrix}=A. \]
Thus $A=QR$, as expected.
\end{example}

The QR factorization will be the major tool used in the proof of the Cholesky
factorization (\ref{ax:7.63}) in the next subsection. For another nice
application of the QR factorization, see the proof of Hadamard's inequality
(\ref{ax:9.66}).

If a QR factorization is available, then it can be used to solve a
corresponding system of linear equations without using Gaussian elimination.
Specifically, suppose $A$ is an $n$-by-$n$ square matrix with linearly
independent columns. Suppose that $b\in\F^n$ and we want to solve the equation
$Ax=b$ for $x=(x_1,\dots,x_n)\in\F^n$ (as usual, we are identifying elements of
$\F^n$ with $n$-by-$1$ column vectors).

Suppose $A=QR$, where $Q$ is unitary and $R$ is upper triangular with only
positive numbers on its diagonal ($Q$ and $R$ are computable from $A$ using
just the Gram--Schmidt procedure, as shown in the proof of \ref{ax:7.58}). The
equation $Ax=b$ is equivalent to the equation $QRx=b$. Multiplying both sides of
this last equation by $Q^*$ on the left and using \ref{ax:7.57}(d) gives the
equation
\[ Rx=Q^*b. \]

The matrix $Q^*$ is the conjugate transpose of the matrix $Q$. Thus computing
$Q^*b$ is straightforward. Because $R$ is an upper-triangular matrix with
positive numbers on its diagonal, the system of linear equations represented by
the equation above can quickly be solved by first solving for $x_n$, then for
$x_{n-1}$, and so on.

\subsection{Cholesky Factorization}

We begin this subsection with a characterization of positive invertible
operators in terms of inner products.

\begin{result}[Positive invertible operator]\label{ax:7.61}
A self-adjoint operator $T\in\Lin(V)$ is a positive invertible operator if and
only if $\ip{Tv}{v}>0$ for every nonzero $v\in V$.
\end{result}

\begin{proof}
First suppose $T$ is a positive invertible operator. If $v\in V$ and $v\ne0$,
then because $T$ is invertible we have $Tv\ne0$. This implies that
$\ip{Tv}{v}\ne0$ (by \ref{ax:7.43}). Hence $\ip{Tv}{v}>0$.

To prove the implication in the other direction, suppose now that
$\ip{Tv}{v}>0$ for every nonzero $v\in V$. Thus $Tv\ne0$ for every nonzero
$v\in V$. Hence $T$ is injective. Thus $T$ is invertible, as desired.
\end{proof}

The next definition transfers the result above to the language of matrices.
Here we are using the usual Euclidean inner product on $\F^n$ and identifying
elements of $\F^n$ with $n$-by-$1$ column vectors.

\begin{definition}[Positive definite]\label{ax:7.62}
A matrix $B\in\F^{n,n}$ is called \emph{positive definite} if $B^*=B$ and
\[ \ip{Bx}{x}>0 \]
for every nonzero $x\in\F^n$.
\end{definition}

A matrix is upper triangular if and only if its conjugate transpose is lower
triangular (meaning that all entries above the diagonal are $0$). The
factorization below, which has important consequences in computational linear
algebra, writes a positive definite matrix as the product of a lower triangular
matrix and its conjugate transpose.

Our next result is solely about matrices, although the proof makes use of the
identification of results about operators with results about square matrices.
In the result below, if the matrix $B$ is in $\F^{n,n}$, then the matrix $R$ is
also in $\F^{n,n}$.

\begin{result}[Cholesky factorization]\label{ax:7.63}
Suppose $B$ is a positive definite matrix. Then there exists a unique
upper-triangular matrix $R$ with only positive numbers on its diagonal such that
\[ B=R^*R. \]
\end{result}

\begin{proof}
Because $B$ is positive definite, there exists an invertible square matrix $A$
of the same size as $B$ such that $B=A^*A$ [by the equivalence of (a) and (f) in
\ref{ax:7.38}].

Let $A=QR$ be the QR factorization of $A$ (see \ref{ax:7.58}), where $Q$ is
unitary and $R$ is upper triangular with only positive numbers on its diagonal.
Then $A^*=R^*Q^*$.

\marginnote{Andr\'e-Louis Cholesky (1875--1918) discovered this factorization,
which was published posthumously in 1924.}%
Thus
\[ B=A^*A=R^*Q^*QR=R^*R, \]
as desired.

To prove the uniqueness part of this result, suppose $S$ is an upper-triangular
matrix with only positive numbers on its diagonal and $B=S^*S$. The matrix $S$
is invertible because $B$ is invertible (see Exercise~11 in Section~3D).
Multiplying both sides of the equation $B=S^*S$ by $S^{-1}$ on the right gives
the equation $BS^{-1}=S^*$.

Let $A$ be the matrix from the first paragraph of this proof. Then
\begin{align*}
(AS^{-1})^*(AS^{-1}) &= (S^*)^{-1}A^*AS^{-1}\\
  &= (S^*)^{-1}BS^{-1}\\
  &= (S^*)^{-1}S^*\\
  &= I.
\end{align*}
Thus $AS^{-1}$ is unitary.

Hence $A=(AS^{-1})S$ is a factorization of $A$ as the product of a unitary
matrix and an upper-triangular matrix with only positive numbers on its
diagonal. The uniqueness of the QR factorization, as stated in \ref{ax:7.58},
now implies that $S=R$.
\end{proof}

In the first paragraph of the proof above, we could have chosen $A$ to be the
unique positive definite matrix that is a square root of $B$ (see
\ref{ax:7.39}). However, the proof was presented with the more general choice of
$A$ because for specific positive definite matrices $B$, it may be easier to
find a different choice of $A$.

% ------------------------------------------------------------------ exercises
\exercisesheading{7D}

\begin{exercise}
Suppose $\dim V\ge2$ and $S\in\Lin(V,W)$. Prove that $S$ is an isometry if and
only if $Se_1,Se_2$ is an orthonormal list in $W$ for every orthonormal list
$e_1,e_2$ of length two in $V$.
\end{exercise}

\begin{exercise}
Suppose $T\in\Lin(V,W)$ and $T\ne0$. Prove that $T$ is a scalar multiple of an
isometry if and only if $T$ preserves orthogonality.
\exnote{The phrase ``$T$ preserves orthogonality'' means that
$\ip{Tu}{Tv}=0$ for all $u,v\in V$ such that $\ip{u}{v}=0$.}
\end{exercise}

\begin{exercise}\leavevmode
\begin{parts}
\item Show that the product of two unitary operators on $V$ is a unitary
  operator.
\item Show that the inverse of a unitary operator on $V$ is a unitary operator.
\end{parts}
\exnote{This exercise shows that the set of unitary operators on $V$ is a
group, where the group operation is the usual product of two operators.}
\end{exercise}

\begin{exercise}
Suppose $\F=\C$ and $A,B\in\Lin(V)$ are self-adjoint. Show that $A+iB$ is
unitary if and only if $AB=BA$ and $A^2+B^2=I$.
\end{exercise}

\begin{exercise}
Suppose $S\in\Lin(V)$. Prove that the following are equivalent.
\begin{parts}
\item $S$ is a self-adjoint unitary operator.
\item $S=2P-I$ for some orthogonal projection $P$ on $V$.
\item There exists a subspace $U$ of $V$ such that $Su=u$ for every $u\in U$
  and $Sw=-w$ for every $w\in U^\perp$.
\end{parts}
\end{exercise}

\begin{exercise}
Suppose $T_1,T_2$ are both normal operators on $\F^3$ with $2,5,7$ as
eigenvalues. Prove that there exists a unitary operator $S\in\Lin(\F^3)$ such
that $T_1=S^*T_2S$.
\end{exercise}

\begin{exercise}
Give an example of two self-adjoint operators $T_1,T_2\in\Lin(\F^4)$ such that
the eigenvalues of both operators are $2,5,7$ but there does not exist a unitary
operator $S\in\Lin(\F^4)$ such that $T_1=S^*T_2S$. Be sure to explain why there
is no unitary operator with the required property.
\end{exercise}

\begin{exercise}
Prove or give a counterexample: If $S\in\Lin(V)$ and there exists an
orthonormal basis $e_1,\dots,e_n$ of $V$ such that $\norm{Se_k}=1$ for each
$e_k$, then $S$ is a unitary operator.
\end{exercise}

\begin{exercise}
Suppose $\F=\C$ and $T\in\Lin(V)$. Suppose every eigenvalue of $T$ has absolute
value $1$ and $\norm{Tv}\le\norm{v}$ for every $v\in V$. Prove that $T$ is a
unitary operator.
\end{exercise}

\begin{exercise}
Suppose $\F=\C$ and $T\in\Lin(V)$ is a self-adjoint operator such that
$\norm{Tv}\le\norm{v}$ for all $v\in V$.
\begin{parts}
\item Show that $I-T^2$ is a positive operator.
\item Show that $T+i\sqrt{I-T^2}$ is a unitary operator.
\end{parts}
\end{exercise}

\begin{exercise}
Suppose $S\in\Lin(V)$. Prove that $S$ is a unitary operator if and only if
\[ \{Sv : v\in V \text{ and } \norm{v}\le1\}=\{v\in V : \norm{v}\le1\}. \]
\end{exercise}

\begin{exercise}
Prove or give a counterexample: If $S\in\Lin(V)$ is invertible and
$\norm{S^{-1}v}=\norm{Sv}$ for every $v\in V$, then $S$ is unitary.
\end{exercise}

\begin{exercise}
Explain why the columns of a square matrix of complex numbers form an
orthonormal list in $\C^n$ if and only if the rows of the matrix form an
orthonormal list in $\C^n$.
\end{exercise}

\begin{exercise}
Suppose $v\in V$ with $\norm{v}=1$ and $b\in\F$. Also suppose $\dim V\ge2$.
Prove that there exists a unitary operator $S\in\Lin(V)$ such that
$\ip{Sv}{v}=b$ if and only if $\abs{b}\le1$.
\end{exercise}

\begin{exercise}
Suppose $T$ is a unitary operator on $V$ such that $T-I$ is invertible.
\begin{parts}
\item Prove that $(T+I)(T-I)^{-1}$ is a skew operator (meaning that it equals
  the negative of its adjoint).
\item Prove that if $\F=\C$, then $i(T+I)(T-I)^{-1}$ is a self-adjoint
  operator.
\end{parts}
\exnote{The function $z\mapsto i(z+1)(z-1)^{-1}$ maps the unit circle in $\C$
(except for the point $1$) to $\R$. Thus (b) illustrates the analogy between
the unitary operators and the unit circle in $\C$, along with the analogy
between the self-adjoint operators and $\R$.}
\end{exercise}

\begin{exercise}
Suppose $\F=\C$ and $T\in\Lin(V)$ is self-adjoint. Prove that
$(T+iI)(T-iI)^{-1}$ is a unitary operator and $1$ is not an eigenvalue of this
operator.
\end{exercise}

\begin{exercise}
Explain why the characterizations of unitary matrices given by \ref{ax:7.57}
hold.
\end{exercise}

\begin{exercise}
A square matrix $A$ is called \emph{symmetric} if it equals its transpose. Prove
that if $A$ is a symmetric matrix with real entries, then there exists a unitary
matrix $Q$ with real entries such that $Q^*AQ$ is a diagonal matrix.
\end{exercise}

\begin{exercise}
Suppose $n$ is a positive integer. For this exercise, we adopt the notation that
a typical element $z$ of $\C^n$ is denoted by $z=(z_0,z_1,\dots,z_{n-1})$.
Define linear functionals $\omega_0,\omega_1,\dots,\omega_{n-1}$ on $\C^n$ by
\[ \omega_j(z_0,z_1,\dots,z_{n-1})
   =\frac{1}{\sqrt{n}}\sum_{m=0}^{n-1}z_me^{-2\pi ijm/n}. \]
The \emph{discrete Fourier transform} is the operator
$\mathcal{F}\colon\C^n\to\C^n$ defined by
\[ \mathcal{F}z=\bigl(\omega_0(z),\omega_1(z),\dots,\omega_{n-1}(z)\bigr). \]
\begin{parts}
\item Show that $\mathcal{F}$ is a unitary operator on $\C^n$.
\item Show that if $(z_0,\dots,z_{n-1})\in\C^n$ and $z_n$ is defined to equal
  $z_0$, then
  \[ \mathcal{F}^{-1}(z_0,z_1,\dots,z_{n-1})=\mathcal{F}(z_n,z_{n-1},\dots,z_1). \]
\item Show that $\mathcal{F}^4=I$.
\end{parts}
\exnote{The discrete Fourier transform has many important applications in data
analysis. The usual Fourier transform involves expressions of the form
$\int_{-\infty}^{\infty}f(x)e^{-2\pi itx}\,dx$ for complex-valued integrable
functions $f$ defined on $\R$.}
\end{exercise}

\begin{exercise}
Suppose $A$ is a square matrix with linearly independent columns. Prove that
there exist unique matrices $R$ and $Q$ such that $R$ is lower triangular with
only positive numbers on its diagonal, $Q$ is unitary, and $A=RQ$.
\end{exercise}
